\documentclass[11pt,a4paper]{article}
\usepackage{gtu-el-yazisi}
\begin{document}
\notbilgi{Sayısal İşaret İşleme}{Sayısal İşaret İşlemenin Temelleri}{Dsp-ders-1.pdf}

\sayfa{1}
\section{Sayısal İşaret İşlemenin Temelleri}
(Discrete-Time Signal Processing) \textit{(Kitap adı)}

\begin{itemize}
  \item 2 tane kısa sınav \quad 20 dk
  \item Kısa sınav \%10
  \item Ödevler \%5 \quad Her ödev 1 puan
  \item Ara sınav \%35
  \item Final sınavı \%50
\end{itemize}

\elcizimi[0.6]{sekiller/Dsp-ders-1-s1-1.png}[1]

\elcizimi[0.6]{sekiller/Dsp-ders-1-s1-2.png}[1]

Kuantalarsak dijital sinyal olur.

Biz ayrık zamanla çalışacağız.

\elcizimi[0.4]{sekiller/Dsp-ders-1-s1-3.png}[1]

\[ x(t) \xrightarrow[t=nT]{\text{sampling}} x[n] \]
\[ x[n] = x(t)\Big|_{t=nT} \]

$T$: period (sec)

$\dfrac{1}{T} = f_s$ : sampling frequency [Hz]

\sayfa{2}
\[ x[n] = [\quad 1 \quad \underset{\substack{\uparrow\\ n=7}}{10} \quad 11 \quad 17 \quad 0\ \ 0\ \ 0]\ 0\ \ 0 \]

\section{Önemli işaretler}

\[ \delta[n] = \begin{cases} 1 & n=0 \\ 0 & \text{o.w.} \end{cases} \]

\[ u[n] = \begin{cases} 1 & n \ge 0 \\ 0 & \text{o.w.} \end{cases} \]

\[ x[n] = a^n \quad \Longrightarrow \quad x[n] = r^n e^{j\theta n} \]

Genlik: $|x[n]| = |r^n|$

$\angle x[n] = \theta n$ radyan

a complex olsun \quad $\overset{\Uparrow}{a = re^{j\theta}}$

\[ x[n] = r^n e^{j\theta n} = \underbrace{r^n \cos\theta n}_{\text{reel}} + j\,\underbrace{r^n \sin\theta n}_{\text{imajiner}} \]

\elcizimi[0.65]{sekiller/Dsp-ders-1-s2-1.png}[2]

\elcizimi[0.62]{sekiller/Dsp-ders-1-s2-2.png}[2]

\begin{align*}
x[n] &= \cdots + x[-1]\,\delta[n+1] + x[0]\,\delta[n] + x[1]\,\delta[n-1] + \\
&\quad + x[2]\,\delta[n-2] + \cdots + x[k]\,\delta[n-k]
\end{align*}

\[ x[n] = \sum_{k=-\infty}^{\infty} x[k]\,\delta[n-k] \]

$\rightarrow$ LTI bir sistemin girişini bu şekilde bilirsek, bir girişin diğerini bilirsek, her giriş için çıkışı biliriz.

\elcizimi[0.45]{sekiller/Dsp-ders-1-s2-3.png}[2]

\sayfa{3}
\elcizimi[0.42]{sekiller/Dsp-ders-1-s3-1.png}[3]

\[ x[n] = \underset{\substack{\uparrow\\ \text{Bu işaret}\\ \text{periyodiktir.}}}{A e^{j\omega_0 n}} = A e^{j\omega_0 n}\,(e^{j2\pi})^n = \underset{\substack{\downarrow\\ \text{$2\pi$ de periyodik}}}{A e^{j(\omega_0+2\pi)n}} \]

\elcizimi[0.38]{sekiller/Dsp-ders-1-s3-2.png}[3]

\begin{align*}
x[n] &= A\cos(\omega_0 n) = A\cos(\omega_0 n + 2\pi k) \\
&= A\cos(\omega_0 (n+N)) = A\cos(\omega_0 n + \omega_0 N)
\end{align*}

\[ x[n+N] = x[n] \ \text{ for all } n \ \Longrightarrow \ \text{işaret periyodiktir} \qquad\qquad 2\pi k = \omega_0 N \]

\[ x_1[n] = \cos\Big(\underbrace{\frac{\pi}{4}}_{\omega_0}\, n\Big) \]

\[ \omega_0 = \frac{\pi}{4} \qquad \text{şartımız} \qquad 2\pi\,\overset{\substack{\text{tamsayı}\\ \text{(0 hariç)}\\ \downarrow}}{k} = \frac{\pi}{4}\,\overset{\substack{\text{bu da tamsayı}\\ \text{olmak zorunda}\\ \downarrow}}{N} \]

\[ 8k = N \]

$k=1$ için $N=\underset{\text{(sample)}}{8}$ çıkar

Her 8 örnekte bir kendini tekrarlar.

2 tam periyot için $n=0$'dan $n=15$'e kadar çizdirilir.

\[ x_2[n] = \cos\Big(\frac{3\pi}{8}\,n\Big) \]

\[ \omega_0 = \frac{3\pi}{8} \qquad\qquad 2\pi k = \frac{3\pi}{8}\,N \]

\[ \frac{16k}{3} = N \]

$k=3$ için $\kutula{N=16\ \mathrm{sample}}$

\[ x_3[n] = \cos(3n) \]

\[ \omega_0 = 3 \qquad\qquad 2\pi k = 3N \qquad\qquad \frac{2\pi}{3}\,k = N \qquad \text{periyodik değil} \]

$k$ ve $N$ aynı zamanda tamsayı olamıyor.

\sayfa{4}
$x_1[n]$ \quad $N_1$

$x_2[n]$ \quad $N_2$

$x_3[n] = x_1[n] + x_2[n]$ \quad periyodik \quad ekok($N_1,N_2$) ile

$\star$ \quad $x_1[n]\cdot x_2[n]$ \quad periyodik mı?? \quad ödev

\[ x[n] \Longrightarrow \text{Enerjisi} \quad \sum_{k=-\infty}^{\infty} |x[k]|^2 \]

\[ \text{Ortalama güç} \quad \lim_{K\to\infty} \frac{1}{2K+1} \sum_{k=-K}^{K} |x[k]|^2 \]

Eğer bir işaretin enerjisi $E_x < \infty$ ise ona enerji işareti denir ve ortalama gücü sıfırdır.

Güç işaretlerinin gücü $0$'dan farklı, enerjisi sonsuz.

\begin{align*}
x_a(t) &= \cos(\Omega_0 t) = \cos(2\pi F_0 t) \\
x_a(t)\Big|_{t=nT} &= x[n] = \cos(\Omega_0 nT) = \cos\Big(\overset{\substack{\text{Frekansı normalize olur}\\ \downarrow}}{\frac{\Omega_0}{F_s}}\, n\Big) \\
&= \cos\Big(2\pi\frac{F_0}{F_s}\,n\Big) \qquad\qquad T = \frac{1}{F_s}
\end{align*}

Örnekleme frekansına göre normalize edilen hali:
\[ {}= \cos(\omega_0 n) = \cos(2\pi f_0 n) \]

\elcizimi[0.9]{sekiller/Dsp-ders-1-s4-1.png}[4]

\sayfa{5}
\[ \cos(\underbrace{\omega_0}_{\text{radyan}}\, n) = \cos(\alpha) \qquad\qquad \text{? Örnek = radyan} \]

\[ \cos\big(\underbrace{2\pi}_{\text{rad}}\, f\, \underset{\uparrow}{n}\big) \]

$\rightarrow$ $\cos(0{,}1\pi n)$ işaretini çizdirin

\[
\begin{gathered}
\cos(0{,}2\pi n) \\ \vdots \\ \cos(0{,}9\pi n) \\ \cos(\pi n) \\ \cos(1{,}1\pi n)
\end{gathered}
\]

\end{document}
